%%
%% MScheme R4RS Conformance Report
%%
\documentclass[12pt]{report}
\usepackage[margin=1in]{geometry}
\usepackage{longtable}
\usepackage{booktabs}
\usepackage{listings}
\usepackage{xcolor}
\usepackage{hyperref}
\usepackage{textcomp}
\usepackage{enumitem}

\lstdefinelanguage{Scheme}{
  morekeywords={define,lambda,if,cond,else,let,let*,letrec,begin,
    set!,quote,quasiquote,unquote,unquote-splicing,and,or,not,
    case,do,delay,force,map,for-each,apply,call-with-current-continuation,
    call/cc,require-modules,load,macro,define-global-symbol,
    import,module,library,program,SchemeStubs,Smodule,ExportSchemeStubs,
    importSchemeStubs},
  sensitive=true,
  morecomment=[l]{;},
  morecomment=[l]{\%},
  morestring=[b]",
  literate={lambda}{$\lambda$}1,
}

\lstdefinelanguage{Modula3}{
  morekeywords={MODULE,INTERFACE,IMPORT,FROM,PROCEDURE,BEGIN,END,
    VAR,CONST,TYPE,RECORD,OBJECT,METHODS,OVERRIDES,RETURN,
    IF,THEN,ELSE,ELSIF,WHILE,DO,FOR,TO,BY,LOOP,EXIT,REPEAT,UNTIL,
    CASE,OF,WITH,TRY,EXCEPT,FINALLY,RAISE,RAISES,
    NEW,ARRAY,REF,REFANY,SET,BOOLEAN,INTEGER,REAL,LONGREAL,CHAR,TEXT,
    TRUE,FALSE,NIL,BRANDED,REVEAL,EVAL,NOT,AND,OR,IN,
    ISTYPE,NARROW,TYPECODE},
  sensitive=true,
  morecomment=[n]{(*}{*)},
  morestring=[b]",
}

\lstset{
  basicstyle=\ttfamily\small,
  breaklines=true,
  frame=single,
  backgroundcolor=\color{gray!5},
  keywordstyle=\bfseries\color{blue!70!black},
  commentstyle=\itshape\color{green!50!black},
  stringstyle=\color{red!60!black},
  showstringspaces=false,
  columns=flexible,
  tabsize=2,
  xleftmargin=2em,
  framexleftmargin=1.5em,
}

\title{\bfseries\sffamily\fontsize{24.88}{30}\selectfont
  MScheme R4RS Conformance Report}
\author{Mika Nystr\"om \\ {\tt mika@alum.mit.edu}}
\date{March 2026}

\begin{document}
\maketitle
\thispagestyle{empty}

\newpage
\tableofcontents
\newpage

%======================================================================
\chapter{Introduction}
%======================================================================

This report is a systematic audit of MScheme against the
\emph{Revised$^4$ Report on the Algorithmic Language Scheme
(R4RS)}~\cite{r4rs}.  MScheme is a Scheme interpreter written in
Modula-3~\cite{spwm3}, descended from Peter Norvig's
JScheme~\cite{jscheme}.

The audit is based on source code review of the interpreter
({\tt mscheme/src/}) and the Scheme compiler
({\tt schemecompiler/src/}).  Each section of R4RS is examined
in turn for:

\begin{enumerate}
\item \textbf{Completeness} --- which required and optional features
  are present, and which are missing.
\item \textbf{Implementation-defined behavior} --- where R4RS grants
  freedom to the implementation and how MScheme exercises that freedom.
\item \textbf{Deviations} --- cases where MScheme's behavior differs
  from R4RS requirements.
\end{enumerate}

Section~\ref{ch:extensions} catalogs MScheme extensions beyond the
scope of R4RS.

%======================================================================
\chapter{Lexical Conventions (R4RS \S2)}
%======================================================================

\section{Identifiers}

MScheme identifiers follow R4RS conventions.  The extended alphabetic
characters ({\tt !\ \$\ \%\ \&\ *\ /\ :\ <\ =\ >\ ?\ \textasciitilde\
\textasciicircum\ \_}) are accepted as part of identifiers.

Identifiers are case-insensitive: the reader converts all symbols to
lower case (conformant with R4RS~\S2).

\section{Comments}

Line comments beginning with {\tt ;} are supported (conformant).

\section{String Literals}

The escape sequences {\tt \textbackslash"} and {\tt \textbackslash\textbackslash}
work correctly.  The sequence {\tt \textbackslash n} in a string literal produces
the literal character~{\tt n}, \emph{not} a newline.  This is conformant:
R4RS~\S2 only requires {\tt \textbackslash"} and {\tt \textbackslash\textbackslash}.

\section{Number Literal Prefixes}

\begin{itemize}
\item {\tt \#d} (decimal): silently accepted (no-op, since decimal is the default).
\item {\tt \#e} (exact) and {\tt \#i} (inexact): silently ignored.  Since
  MScheme has no exact/inexact distinction, this is harmless.
\item {\tt \#b} (binary), {\tt \#o} (octal), {\tt \#x} (hexadecimal):
  emit a warning ``not implemented'' and fall through to decimal parsing.
\end{itemize}

\noindent
The {\tt string->number} procedure, by contrast, fully supports radix
2--16 via its optional second argument.

\section{Character Literals}

{\tt \#\textbackslash\emph{c}}, {\tt \#\textbackslash space}, and
{\tt \#\textbackslash newline} are all supported (conformant).

\section{Vector Literals}

{\tt \#(\ldots)} vector literals are supported (conformant).

%======================================================================
\chapter{Expressions and Special Forms (R4RS \S4)}
%======================================================================

\section{Primitive Expression Types}

All R4RS primitive expression types are present:

\begin{center}
\begin{tabular}{ll}
\toprule
\textbf{Form} & \textbf{Status} \\
\midrule
Variable reference & Present \\
{\tt (quote \emph{datum})} & Present \\
{\tt (lambda \emph{formals} \emph{body})} & Present (incl.\ rest args) \\
{\tt (if \emph{test} \emph{consequent} \emph{alternate})} & Present \\
{\tt (set!\ \emph{variable} \emph{expression})} & Present \\
Procedure call & Present \\
\bottomrule
\end{tabular}
\end{center}

\section{Derived Expression Types}

All R4RS derived expression types are implemented as macros in
{\tt SchemePrimitives.i3}:

\begin{center}
\begin{tabular}{ll}
\toprule
\textbf{Form} & \textbf{Status} \\
\midrule
{\tt cond} & Present (incl.\ {\tt =>} syntax) \\
{\tt case} & Present (see deviation below) \\
{\tt and}, {\tt or} & Present \\
{\tt let}, {\tt let*}, {\tt letrec} & Present \\
Named {\tt let} & Present \\
{\tt begin} & Present \\
{\tt do} & Present \\
{\tt delay} & Present \\
{\tt quasiquote} & Present \\
\bottomrule
\end{tabular}
\end{center}

\paragraph{Deviation: {\tt case} macro.}
The {\tt case} macro expansion uses {\tt member} (which calls
{\tt equal?}) to test clause membership.  R4RS~\S4.2.1 specifies
that the keys are compared using {\tt eqv?}, so the correct
procedure would be {\tt memv}.  In practice this only matters for
number comparisons where {\tt equal?} and {\tt eqv?} might differ
(e.g., with exact/inexact distinctions), which MScheme does not
have, so the deviation is benign.

\paragraph{{\tt cond} with {\tt =>}.}
The {\tt =>} syntax in {\tt cond} clauses is handled directly in the
evaluator ({\tt Scheme.m3}, {\tt ReduceCond}), not via macro
expansion.

\section{Internal Definitions}

Internal {\tt define} forms at the beginning of a lambda body are
supported.  They are handled at runtime by the evaluator.
{\tt SchemeClosure.Init} recognizes internal definitions but does not
convert them to an equivalent {\tt letrec}; instead, it disables its
body optimization pass and lets the defines fall through to the
normal evaluation path (conformant in behavior with R4RS~\S5.2.2,
though not via the transformation suggested by the report).

\section{Tail Recursion}

MScheme implements proper tail recursion via a {\tt LOOP} in the main
evaluator.  The following positions are recognized as tail positions
(conformant with R4RS~\S3.5):

\begin{itemize}
\item Last expression in a {\tt lambda} body.
\item Both branches of {\tt if}.
\item Last selected branch of {\tt cond}.
\item Last expression in {\tt and} / {\tt or}.
\item Last expression in {\tt begin}, {\tt let}, {\tt let*}, {\tt letrec}.
\item Body of named {\tt let}.
\item Last expression in the {\tt when} clause of {\tt do}.
\end{itemize}

%======================================================================
\chapter{Standard Procedures (R4RS \S6)}
%======================================================================

This chapter reviews each subsection of R4RS~\S6.

%----------------------------------------------------------------------
\section{Booleans (\S6.1)}
%----------------------------------------------------------------------

{\tt not} and {\tt boolean?} are present.  Complete.

%----------------------------------------------------------------------
\section{Equivalence Predicates (\S6.2)}
%----------------------------------------------------------------------

{\tt eq?}, {\tt eqv?}, and {\tt equal?} are all present.

\paragraph{Note.} {\tt eqv?} compares {\tt LONGREAL} values directly.
Since MScheme has no exact/inexact distinction, {\tt (eqv?\ 1\ 1.0)}
is {\tt \#t}.  See Section~\ref{sec:impl-defined} for further
discussion.

%----------------------------------------------------------------------
\section{Pairs and Lists (\S6.3)}
%----------------------------------------------------------------------

All R4RS pair and list procedures are present:

\begin{itemize}[nosep]
\item {\tt cons}, {\tt car}, {\tt cdr}, {\tt set-car!}, {\tt set-cdr!}
\item All 28 {\tt c\ldots r} compositions ({\tt caar} through {\tt cddddr})
\item {\tt null?}, {\tt list?}, {\tt list}, {\tt length}, {\tt append},
  {\tt reverse}
\item {\tt list-tail}, {\tt list-ref}
\item {\tt memq}, {\tt memv}, {\tt member}
\item {\tt assq}, {\tt assv}, {\tt assoc}
\end{itemize}

\noindent Complete.

%----------------------------------------------------------------------
\section{Symbols (\S6.4)}
%----------------------------------------------------------------------

{\tt symbol?}, {\tt symbol->string}, {\tt string->symbol} are present.
Complete.

%----------------------------------------------------------------------
\section{Numbers (\S6.5)}
\label{sec:numbers}
%----------------------------------------------------------------------

This is the area of greatest divergence from R4RS.

\subsection{Number Tower}

MScheme has a single numeric type: IEEE~754 double precision
(Modula-3 {\tt LONGREAL}).  There are no exact integers, no rationals,
and no complex numbers.  This means:

\begin{itemize}
\item {\tt complex?} = {\tt real?} = {\tt number?} (all return {\tt \#t}
  for any number).
\item {\tt exact?} = {\tt integer?} = {\tt rational?} (true when the
  double value is integer-valued).
\item {\tt inexact?} returns {\tt \#t} when the value is \emph{not}
  integer-valued.
\end{itemize}

R4RS acknowledges that an implementation ``may use fixnums, flonums,
or perhaps some other representation'' but expects at least an
exact/inexact distinction.  MScheme collapses this distinction
entirely.

\subsection{Arithmetic and Comparisons}

All R4RS arithmetic operations are present: {\tt +}, {\tt -}, {\tt *},
{\tt /}, {\tt =}, {\tt <}, {\tt >}, {\tt <=}, {\tt >=}.

Multi-argument forms work correctly.  Division by zero
signals an error.

\subsection{Rounding, GCD, LCM}

{\tt floor}, {\tt ceiling}, {\tt truncate}, {\tt round},
{\tt gcd}, {\tt lcm} are all present.

\subsection{Transcendentals}

{\tt exp}, {\tt log}, {\tt sin}, {\tt cos}, {\tt tan}, {\tt asin},
{\tt acos}, {\tt atan}, {\tt sqrt}, {\tt expt} are present.

\paragraph{Deviation.}
{\tt atan} accepts only one argument.  The two-argument form
{\tt (atan \emph{y} \emph{x})} required by R4RS is missing.

\subsection{Number I/O}

{\tt number->string} and {\tt string->number} are present.  Both
support the optional radix argument with bases 2--16.

\subsection{Missing Procedures}

\begin{center}
\begin{tabular}{lll}
\toprule
\textbf{Procedure} & \textbf{R4RS Status} & \textbf{Notes} \\
\midrule
{\tt exact->inexact} & Essential & Missing \\
{\tt inexact->exact} & Essential & Missing \\
{\tt numerator} & Non-essential & Missing \\
{\tt denominator} & Non-essential & Missing \\
{\tt rationalize} & Non-essential & Missing \\
{\tt make-rectangular} & Non-essential & No complex numbers \\
{\tt make-polar} & Non-essential & No complex numbers \\
{\tt real-part} & Non-essential & No complex numbers \\
{\tt imag-part} & Non-essential & No complex numbers \\
{\tt magnitude} & Non-essential & No complex numbers \\
{\tt angle} & Non-essential & No complex numbers \\
2-arg {\tt atan} & Essential & Missing \\
\bottomrule
\end{tabular}
\end{center}

\noindent
{\tt exact->inexact} and {\tt inexact->exact} are marked essential in
R4RS.  They could be trivially implemented as identity functions given
MScheme's single-type number representation.  The two-argument form of
{\tt atan} could be implemented via {\tt Math.atan2}.

%----------------------------------------------------------------------
\section{Characters (\S6.6)}
%----------------------------------------------------------------------

All R4RS character procedures are present:

\begin{itemize}[nosep]
\item {\tt char?}
\item Comparisons: {\tt char=?}, {\tt char<?}, {\tt char>?},
  {\tt char<=?}, {\tt char>=?}
\item Case-insensitive comparisons: {\tt char-ci=?}, {\tt char-ci<?},
  {\tt char-ci>?}, {\tt char-ci<=?}, {\tt char-ci>=?}
\item {\tt char-alphabetic?}, {\tt char-numeric?}, {\tt char-whitespace?},
  {\tt char-upper-case?}, {\tt char-lower-case?}
\item {\tt char->integer}, {\tt integer->char}
\item {\tt char-upcase}, {\tt char-downcase}
\end{itemize}

\noindent
Complete.  Character ordering is based on the Modula-3 {\tt CHAR}
ordinal (ISO~Latin-1, a superset of ASCII).

%----------------------------------------------------------------------
\section{Strings (\S6.7)}
%----------------------------------------------------------------------

Present:
\begin{itemize}[nosep]
\item {\tt string?}, {\tt make-string}, {\tt string}
\item {\tt string-length}, {\tt string-ref}, {\tt string-set!}
\item Comparisons: {\tt string=?}, {\tt string<?}, {\tt string>?},
  {\tt string<=?}, {\tt string>=?}
\item Case-insensitive comparisons: {\tt string-ci=?}, {\tt string-ci<?},
  {\tt string-ci>?}, {\tt string-ci<=?}, {\tt string-ci>=?}
\item {\tt substring}, {\tt string-append}
\item {\tt string->list}, {\tt list->string}
\end{itemize}

\paragraph{Missing (non-essential).}
{\tt string-copy} and {\tt string-fill!}.

%----------------------------------------------------------------------
\section{Vectors (\S6.8)}
%----------------------------------------------------------------------

Present:
\begin{itemize}[nosep]
\item {\tt vector?}, {\tt make-vector}, {\tt vector}
\item {\tt vector-length}, {\tt vector-ref}, {\tt vector-set!}
\item {\tt vector->list}, {\tt list->vector}
\end{itemize}

\paragraph{Missing (non-essential).}
{\tt vector-fill!}.

%----------------------------------------------------------------------
\section{Control Features (\S6.9)}
%----------------------------------------------------------------------

{\tt procedure?}, {\tt apply}, {\tt map}, {\tt for-each}, {\tt force}
are all present.

\subsection{{\tt call-with-current-continuation}}

{\tt call-with-current-continuation} (also available as {\tt call/cc})
is present but with significant limitations:

\begin{enumerate}
\item \textbf{Escape-only.}  The captured continuation can only be used
  to exit the {\tt call/cc} form.  It cannot be invoked after the
  {\tt call/cc} expression has returned.  The continuation is
  implemented via Modula-3 exception unwinding.

\item \textbf{Fixed marker.}  The implementation uses a fixed string
  marker {\tt "CallCC123"} to identify the exception frame.  This
  means that nested {\tt call/cc} invocations share the same marker
  and may not behave correctly.

\item \textbf{No re-entry.}  Continuations cannot be stored and later
  invoked to re-enter a computation.  This rules out coroutines,
  cooperative multitasking, and other advanced control flow patterns
  that rely on full continuations.
\end{enumerate}

\noindent
These limitations are inherited from JScheme.  R4RS requires full
continuations, so this is a known deviation.  In practice, escape
continuations suffice for most error-handling and early-return use
cases.

%----------------------------------------------------------------------
\section{Input and Output (\S6.10)}
%----------------------------------------------------------------------

\subsection{Present Procedures}

\begin{itemize}[nosep]
\item Ports: {\tt input-port?}, {\tt output-port?},
  {\tt current-input-port}, {\tt current-output-port}
\item File ports: {\tt open-input-file}, {\tt open-output-file},
  {\tt close-input-port}, {\tt close-output-port}
\item File context: {\tt call-with-input-file}, {\tt call-with-output-file}
\item Reading: {\tt read}, {\tt read-char}, {\tt peek-char},
  {\tt eof-object?}
\item Writing: {\tt write}, {\tt display}, {\tt newline}, {\tt write-char}
\item Loading: {\tt load}
\end{itemize}

\paragraph{Note.}
{\tt write-char} is mapped to the {\tt display} primitive.  This is
correct behavior for characters: {\tt display} writes a character
without the {\tt \#\textbackslash} prefix.

\subsection{Missing (Non-essential)}

\begin{itemize}[nosep]
\item {\tt with-input-from-file}, {\tt with-output-to-file}
\item {\tt char-ready?}
\item {\tt transcript-on}, {\tt transcript-off}
\end{itemize}

%======================================================================
\chapter{Implementation-Defined Behavior}
\label{sec:impl-defined}
%======================================================================

R4RS leaves many details to the implementation.  This section
documents MScheme's choices, organized as a numbered list with R4RS
section references.

\begin{enumerate}[label=\textbf{\arabic*.}, leftmargin=3em]

\item \textbf{Empty list truthiness} (\S3.2).
  R4RS permits {\tt '()} to count as false.  In MScheme, only {\tt \#f}
  is false; {\tt '()} is true.
  \emph{Source:} {\tt SchemeBoolean.m3}: {\tt TruthO(x) = (x \# LFalse)}.

\item \textbf{Number representation} (\S6.5).
  All numbers are IEEE~754 {\tt LONGREAL} (double precision).  There is
  no exact/inexact distinction at the type level.

\item \textbf{{\tt exact?}\ semantics} (\S6.5.5).
  Returns {\tt \#t} when the double value is integer-valued, specifically
  when {\tt d = FLOAT(ROUND(d))} and {\tt ABS(d) < 10\textsuperscript{17}}.
  Functionally equivalent to {\tt integer?}.

\item \textbf{{\tt eqv?}\ on numbers} (\S6.2).
  Compares {\tt LONGREAL} values directly.  Since there is no
  exact/inexact distinction, {\tt (eqv?\ 1\ 1.0)} evaluates to
  {\tt \#t}.  R4RS says {\tt eqv?}\ on numbers of different exactness
  returns {\tt \#f}.

\item \textbf{{\tt set!}\ return value} (\S4.1.6).
  Returns the assigned value.  R4RS says the result is ``unspecified.''

\item \textbf{{\tt define} return value} (\S5.2).
  Returns the symbol being defined.  R4RS says the result is ``unspecified.''

\item \textbf{{\tt if} without alternate} (\S4.1.5).
  Returns {\tt '()} (the empty list) when the test is false.  R4RS says
  the result is ``unspecified.''

\item \textbf{Character ordering} (\S6.6).
  Based on the Modula-3 {\tt CHAR} ordinal, which is ISO~Latin-1
  (a superset of ASCII).

\item \textbf{{\tt eqv?}\ on characters} (\S6.2).
  Characters are represented as shared singletons (an array of 256
  pre-allocated character objects).  Characters with the same code point
  are always {\tt eqv?}.

\item \textbf{{\tt eqv?}\ on empty strings and vectors} (\S6.2).
  Identity comparison.  Two separately allocated empty strings (or
  empty vectors) are \emph{not} {\tt eqv?}.

\item \textbf{{\tt map} evaluation order} (\S6.9).
  Elements are processed left to right.  This is an implementation
  artifact; R4RS does not specify the order.

\item \textbf{Tail positions} (\S4).
  Proper tail recursion is implemented.  The tail positions recognized
  are listed in Section~3.4 of this report.

\item \textbf{{\tt call-with-current-continuation}} (\S6.9).
  Escape-only continuations via Modula-3 exception unwinding.  Fixed
  marker {\tt "CallCC123"}.  Cannot be invoked after the {\tt call/cc}
  has returned.  Nested {\tt call/cc} has issues with the shared marker.
  See Section~4.10.1 for details.

\item \textbf{{\tt write-char}} (\S6.10.3).
  Implemented by delegating to {\tt display}, which is correct for
  characters (writes the character without the {\tt \#\textbackslash}
  prefix).

\item \textbf{{\tt load}} (\S6.10.4).
  Searches {\tt **scheme-load-path**}, checks {\tt SchemeCompiledRegistry}
  for compiled modules, falls back to bundled source.  Returns the filename
  (R4RS says the return value is ``unspecified'').

\end{enumerate}

%======================================================================
\chapter{Gaps Summary}
%======================================================================

Table~\ref{tab:gaps} summarizes all missing or incomplete features.

\begin{table}[ht]
\centering
\caption{Missing and incomplete R4RS features.}
\label{tab:gaps}
\begin{tabular}{lll}
\toprule
\textbf{Procedure / Feature} & \textbf{R4RS Status} & \textbf{Notes} \\
\midrule
\multicolumn{3}{l}{\emph{Essential, missing}} \\
\addlinespace
{\tt exact->inexact}   & Essential & Trivial to add (identity fn) \\
{\tt inexact->exact}   & Essential & Trivial to add (identity fn) \\
2-arg {\tt atan}       & Essential & Use {\tt Math.atan2} \\
\addlinespace
\midrule
\multicolumn{3}{l}{\emph{Non-essential, missing}} \\
\addlinespace
{\tt string-copy}           & Non-essential & --- \\
{\tt string-fill!}          & Non-essential & --- \\
{\tt vector-fill!}          & Non-essential & --- \\
{\tt char-ready?}           & Non-essential & --- \\
{\tt with-input-from-file}  & Non-essential & --- \\
{\tt with-output-to-file}   & Non-essential & --- \\
{\tt transcript-on}         & Non-essential & --- \\
{\tt transcript-off}        & Non-essential & --- \\
{\tt numerator}             & Non-essential & No rationals \\
{\tt denominator}           & Non-essential & No rationals \\
{\tt rationalize}           & Non-essential & No rationals \\
{\tt make-rectangular}      & Non-essential & No complex numbers \\
{\tt make-polar}            & Non-essential & No complex numbers \\
{\tt real-part}             & Non-essential & No complex numbers \\
{\tt imag-part}             & Non-essential & No complex numbers \\
{\tt magnitude}             & Non-essential & No complex numbers \\
{\tt angle}                 & Non-essential & No complex numbers \\
\addlinespace
\midrule
\multicolumn{3}{l}{\emph{Partial / deviant}} \\
\addlinespace
{\tt case} macro            & --- & Uses {\tt equal?} not {\tt eqv?} \\
{\tt \#b}/{\tt \#o}/{\tt \#x} literals & --- & Warning, not implemented \\
{\tt \#e}/{\tt \#i} prefixes  & --- & Silently ignored \\
{\tt call/cc}               & Essential & Escape-only \\
\bottomrule
\end{tabular}
\end{table}

%======================================================================
\chapter{Extensions Beyond R4RS}
\label{ch:extensions}
%======================================================================

MScheme provides a number of features beyond R4RS.
Table~\ref{tab:extensions} catalogs them.

\begin{longtable}{ll}
\caption{MScheme extensions.}
\label{tab:extensions} \\
\toprule
\textbf{Name} & \textbf{Description} \\
\midrule
\endfirsthead
\toprule
\textbf{Name} & \textbf{Description} \\
\midrule
\endhead
\bottomrule
\endfoot
{\tt macro}                 & User-defined syntactic transformers \\
{\tt unwind-protect}        & Exception-safe cleanup \\
{\tt current-environment}   & Access lexical environment \\
{\tt error}                 & Raise error with message \\
{\tt exit}                  & Terminate interpreter \\
{\tt system}                & Shell command execution \\
{\tt random}, {\tt normal}  & Random number generation \\
{\tt time}                  & Benchmarking macro \\
{\tt require-modules}       & Module loading system \\
{\tt eval}                  & Eval with optional environment \\
\addlinespace
\multicolumn{2}{l}{\emph{Mathematics}} \\
\addlinespace
{\tt sinh}, {\tt cosh}, {\tt tanh}    & Hyperbolic trigonometry \\
{\tt asinh}, {\tt acosh}, {\tt atanh}  & Inverse hyperbolic trigonometry \\
\addlinespace
\multicolumn{2}{l}{\emph{Modula-3 interop}} \\
\addlinespace
{\tt number->LONGREAL}      & Convert to Modula-3 {\tt LONGREAL} \\
{\tt number->EXTENDED}      & Convert to Modula-3 {\tt EXTENDED} \\
{\tt number->REAL}          & Convert to Modula-3 {\tt REAL} \\
\addlinespace
\multicolumn{2}{l}{\emph{String/list extensions}} \\
\addlinespace
{\tt string-havesub?}       & Substring search \\
{\tt read-big-int}          & Big integer input \\
\addlinespace
\multicolumn{2}{l}{\emph{Debugging and meta-programming}} \\
\addlinespace
{\tt macro-expand}          & Macro debugging \\
{\tt dump-environment}      & Environment serialization \\
{\tt load-environment!}     & Environment deserialization \\
\addlinespace
\multicolumn{2}{l}{\emph{Performance}} \\
\addlinespace
{\tt eq?-memo}              & Memoization ({\tt eq?} keys) \\
{\tt equal?-memo}           & Memoization ({\tt equal?} keys) \\
{\tt write-noflush}         & Write without flushing \\
{\tt display-noflush}       & Display without flushing \\
\end{longtable}

%======================================================================
\chapter{Conclusion}
%======================================================================

MScheme is a substantially complete R4RS implementation suitable for
practical Scheme programming.  All expression types and special forms
are present.  The standard library covers booleans, equivalence
predicates, pairs and lists, symbols, characters, strings, vectors,
control features, and I/O in full or near-full compliance.

The main gaps lie in three areas:

\begin{enumerate}
\item \textbf{Numeric tower.}  The use of a single {\tt LONGREAL}
  type means there is no exact/inexact distinction, no rational
  arithmetic, and no complex numbers.  Three essential procedures
  ({\tt exact->inexact}, {\tt inexact->exact}, 2-arg {\tt atan}) are
  missing but would be trivial to add.

\item \textbf{Continuations.}  {\tt call-with-current-continuation}
  provides only escape continuations, not full continuations.  This is
  a fundamental limitation of the implementation strategy (Modula-3
  exception unwinding) and is shared with the ancestor JScheme.

\item \textbf{Non-essential I/O and utility procedures.}  A handful of
  non-essential procedures ({\tt string-copy}, {\tt vector-fill!},
  {\tt with-input-from-file}, etc.) are missing.  These are
  straightforward to add.
\end{enumerate}

\noindent
MScheme compensates for these gaps with a rich set of extensions,
including a macro system, module system, Modula-3 interoperation, and
automatic stub generation, making it well suited to its primary role
as an embedded scripting language for Modula-3 applications.

\begin{thebibliography}{99}

\bibitem{r4rs} William Clinger and Jonathan Rees, eds.
Revised$^4$ Report on the Algorithmic Language Scheme.  1991.

\bibitem{spwm3} Greg~Nelson, ed.  {\it Systems Programming with Modula-3.}
Englewood Cliffs, N.J.:\ Prentice-Hall, 1991.

\bibitem{jscheme} Peter Norvig.  JScheme.
{\tt https://norvig.com/jscheme.html}

\end{thebibliography}

\end{document}
